\section{Background}

\subsection{The building and its classrooms}

We counted classrooms from the university's classroom A/V listing for Mundelein Center. Table~\ref{tab:rooms} gives the
count per floor. There are no classrooms on floors 8 and~9, so the one-bank limit at floor~9 splits the
building cleanly into a lower group (floors 1--7, 45 rooms) and an upper group (floors 10--14, 9 rooms) that only
the two-car bank can reach.

\begin{table}[h]
\centering
\caption{Classrooms per floor in Mundelein Center.}
\label{tab:rooms}
\begin{tabular}{lcccccccccccc}
\toprule
Floor & 1 & 2 & 3 & 4 & 5 & 6 & 7 & 10 & 11 & 12 & 13 & 14\\
\midrule
Rooms & 2 & 4 & 5 & 8 & 10 & 11 & 5 & 4 & 1 & 1 & 1 & 2\\
\bottomrule
\end{tabular}
\end{table}

Floor~1 is an atrium lobby, and the stairs from floor~1 to floor~2 feel like two ordinary flights. Students
leave the building through the lobby.

\subsection{The scheduling grid}

The Registrar's standard grid~\cite{luc-grid} defines seven 50-minute Monday--Wednesday--Friday blocks (8:15 to
3:35) and five 75-minute Tuesday--Thursday blocks (8:30 to 3:45) before the evening blocks, which start at
4:15 and are outside the scope of this study. Every consecutive pair of blocks, in either pattern, is separated
by a 15-minute passing period. Our simulation models that period.

Classes seat 30--35 students, but enrollment is often 15--25. Elevator cars are rated for 3{,}000~lb, which is
roughly 17--20 riders by weight; observed loads are four to six, so space, not weight, is the binding limit.
